\subsection{A smaller integral homology gate for \texorpdfstring{$K_{12}$}{K12}}

There is an independent route to the first undecided complete vanishing
degree. It starts from integral Steinberg homology rather than the
automorphic trace polynomial. The following localized form of a known
resolution removes one of the three ranks listed in the earlier
Voronoi computation.

\begin{sourcefact}[K17-ST2-LOCAL]\label{prop:steinberg-two-local}
For every $n\ge3$, put $R_n=\mathbb Z[1/(n+1)!]$ and let
$\mathrm{St}_n=\widetilde H_{n-2}(\mathcal T_n(\mathbb Q),\mathbb Z)$.
Then
\[
 H_2(\mathrm{SL}_n(\mathbb Z);\mathrm{St}_n\otimes R_n)=0,
 \qquad
 H_2(\mathrm{GL}_n(\mathbb Z);\mathrm{St}_n\otimes R_n)=0.
\]
\end{sourcefact}

This is already included in Miller--Patzt--Putman,
Remark~1.18, p.~7 \cite{MillerPatztPutman2025}: the cited
codimension-two proof applies whenever every prime at most $n$
is invertible. Their Lemma~25.3, Claim~2, pp.~56--57, explains
that localization argument for the degree-one model; the
uniform degree-two assertion is stated in Remark~1.18, not in
Claim~2 itself. Our ring $R_n$ satisfies the stated condition.
Moreover their Corollary~E, p.~7, already proves degree-two
vanishing with arbitrary coefficients for $n\ge8$, so the
rank-$11$ case used below holds even over $\mathbb Z$.

Before locating this newer source we independently checked the
weaker $R_n$ form directly from Br\"uck--Miller--Patzt--Sroka--Wilson,
Proposition~8.5 and Lemmas~8.7 and~8.9, pp.~57--59
\cite{BrueckEtAl2024}. That complete recheck remains in the
evidence, with its actual audit events. It is not claimed as a
new result of this run and is not reproved here. The conditional
$K_{12}$ reduction that uses it is given in full next.

\begin{proposition}[K17-K12-TWO-RANKS]\label{prop:k12-two-ranks}
Put $R=\mathbb Z[1/13!]$. If
\[
 H_4(\mathrm{GL}_9(\mathbb Z);\mathrm{St}_9\otimes R)=0,
 \qquad
 H_3(\mathrm{GL}_{10}(\mathbb Z);\mathrm{St}_{10}\otimes R)=0,
                                                               \tag{K23}
\]
then $K_{12}(\mathbb Z)=0$.
Equivalently, it suffices to establish the corresponding localized
vanishings of $H_{12}(\mathrm{Vor}_{\mathrm{GL}_9(\mathbb Z)})$
and $H_{12}(\mathrm{Vor}_{\mathrm{GL}_{10}(\mathbb Z)})$.
Neither of the two hypotheses is established here.
\end{proposition}

\begin{proof}
Use Quillen's rank spectral sequence
\[
 E^1_{a,b}=H_b(\mathrm{GL}_a(\mathbb Z);\mathrm{St}_a\otimes R)
 \ \Longrightarrow\ H_{a+b}(BQ;R).
\]
The exact construction and its Voronoi comparison are recorded in
Dutour Sikiri\'c--Elbaz-Vincent--Kupers--Martinet,
Proposition~1.5 and \S6.1; their Table~5 is on author-preprint p.~17
\cite{DEKM2026}. In total degree $13$, that table and
Theorem~4.5 give zero in ranks $2\le a\le8$ except for
\[
 E^1_{7,6}=R.
\]
Ranks $0$ and $1$ also give zero in this degree: the former is the
point $BQ_0$, and $\mathrm{GL}_1(\mathbb Z)=C_2$ has no positive
homology over $R$. The terms of ranks $9$ and $10$ vanish by (K23). The rank-$11$
term is zero by Source statement~\ref{prop:steinberg-two-local}.
The rank-$12$ term is $H_1(\mathrm{GL}_{12};\mathrm{St}_{12}\otimes R)$,
which vanishes by the localized Church--Putman result; the rank-$13$
term is the zero Steinberg coinvariant group of Lee--Szczarba.
Both precise forms are recorded in Kupers, Lemmas~2.1--2.2 and the
complete localization argument on pp.~2--3
\cite{Kupers2019}. Higher ranks have negative second degree and do
not occur. Thus the only nonzero $E^1$ term on this total-degree
diagonal is the single free rank-one module $R$.

The rational homology of $BQ$ in degree $13$ has rank one.
Indeed, Borel's rational $K$-group calculation and the rational
infinite-loop-space decomposition give an exterior generator in
degree $1$ and polynomial generators in degrees $6,10,14,\ldots$.
The only monomial of degree $13$ is the degree-one generator times
the square of the degree-six generator. These are the same rational
inputs used in the proof of Proposition~6.1 of
\cite{DEKM2026}; the additional degree requires just the
displayed monomial count.

It follows that no differential can have nonzero image in the
surviving rank-one term. More explicitly, inductively that term is
a submodule of $R$ until an incoming nonzero differential occurs.
Any nonzero submodule of a rank-one module over the domain $R$ has
rank one; such an incoming image would kill the only rational
contribution in total degree $13$, contradicting the computed
rational rank. Hence all incoming images are zero. Outgoing
differentials replace the term by a submodule and preserve its
torsionfreeness. The limiting diagonal has only this one term, so
$H_{13}(BQ;R)$ is torsion free.

The Hurewicz map
\[
 K_{12}(\mathbb Z)\otimes R
   =\pi_{13}(BQ)\otimes R\longrightarrow H_{13}(BQ;R)
\]
is injective. For an exact source justification, apply Arlettaz's
stable Hurewicz theorem to the $0$-connected spectrum
$\Sigma\mathbf K(\mathbb Z)$: Theorem~6.6 and Definition~6.5 of
\cite[pp.~40--42]{Arlettaz2000} bound the kernel in degree $13$
by a product of exponents of the first $12$ positive stable stems.
Only primes at most $(12+3)/2$, hence only $2,3,5,7$, divide that
product. The ordinary Hurewicz map of
$BQ=\Omega^\infty\Sigma\mathbf K(\mathbb Z)$ followed by homology
suspension is the stable Hurewicz map, so its kernel is a subgroup
of that same kernel. All these primes are units in $R$.

The finite group $K_{12}(\mathbb Z)\otimes R$ now injects into a
torsion-free module, so it is zero. The primary parts at all primes
inverted in $R$ vanish as well: the $2$-primary part is zero, and
$3,5,7,11,13$ are regular primes. Therefore $K_{12}(\mathbb Z)=0$.
Finally, Proposition~1.5 of \cite{DEKM2026} identifies
$H_q(\mathrm{GL}_n;\mathrm{St}_n)$ with
$H_{q+n-1}(\mathrm{Vor}_{\mathrm{GL}_n})$ after inverting the
indicated finite-stabilizer orders. At $(n,q)=(9,4)$ and $(10,3)$
the Voronoi degree is $12$, proving the equivalent formulation.
\end{proof}

The earlier Voronoi remark listed ranks $9,10,11$ as missing
\cite[Remark~6.6, author-preprint p.~19]{DEKM2026}.
Source statement~\ref{prop:steinberg-two-local} removes rank $11$.
The remaining rank-$9$ and rank-$10$ computations are an alternative
to the complete trace-polynomial certificate, not a conclusion
that the still undecided group already vanishes.
